\hypertarget{c14-functions-and-lambdas}{%
\chapter{C14 FUNCTIONS AND LAMBDAS}\label{c14-functions-and-lambdas}}

\hypertarget{c14-functions-lambdas}{%
\chapter{C++14 FUNCTIONS \& LAMBDAS}\label{c14-functions-lambdas}}

\hypertarget{generic-lambdas}{%
\section{1. Generic Lambdas}\label{generic-lambdas}}

C++14 allows \passthrough{\lstinline!auto!} in lambda parameters,
effectively making them templates.

\begin{lstlisting}[language={C++}]
auto print = [](auto x) {
    std::cout << x << "\n";
};

print(10);      // int
print("hello"); // const char*
\end{lstlisting}

This is shorthand for a struct with a templated
\passthrough{\lstinline!operator()!}.

\hypertarget{generalized-lambda-captures-init-capture}{%
\section{2. Generalized Lambda Captures
(Init-Capture)}\label{generalized-lambda-captures-init-capture}}

C++14 allows initializing variables inside the lambda capture clause.
This is crucial for \textbf{move-only types}.

\begin{lstlisting}[language={C++}]
auto ptr = std::make_unique<int>(10);

// Move ptr into lambda
auto lambda = [p = std::move(ptr)]() {
    std::cout << *p << "\n";
};
// ptr is now nullptr
\end{lstlisting}

\hypertarget{automatic-return-type-deduction}{%
\section{3. Automatic Return Type
Deduction}\label{automatic-return-type-deduction}}

Functions can deduce their return type from the return statement.

\begin{lstlisting}[language={C++}]
auto add(int a, int b) {
    return a + b; // Deduced as int
}

// decltype(auto) preserves references
int& getRef(int& x) { return x; }

decltype(auto) forwardRef(int& x) {
    return getRef(x); // Returns int& (auto would return int)
}
\end{lstlisting}

\begin{center}\rule{0.5\linewidth}{0.5pt}\end{center}
